\section*{5. Results}
\subsection*{5.1 Four-arm cell-policy comparison}
\begin{center}\footnotesize\setlength{\tabcolsep}{3pt}
\begin{tabularx}{\linewidth}{>{\hsize=1.600\hsize\linewidth=\hsize\raggedright\arraybackslash}X>{\hsize=0.800\hsize\linewidth=\hsize\raggedright\arraybackslash}X>{\hsize=0.800\hsize\linewidth=\hsize\raggedright\arraybackslash}X>{\hsize=0.800\hsize\linewidth=\hsize\raggedright\arraybackslash}X}
\toprule
Method & Reverse-ID Top-1 \(\uparrow\) & Sequential AUC3 \(\uparrow\) & CEX AURB3 \(\uparrow\) \\
\midrule
Our harness & 0.803922 & 0.723917 & 0.096127 \\
Our harness + federated evidence & 0.725490 & 0.723470 & 0.096127 \\
Harness-free & 0.764706 & 0.719532 & 0.096127 \\
Cosine baseline & 0.705882 & 0.721687 & N/A \\
\bottomrule
\end{tabularx}
\end{center}
Table 1. Exploratory cell-policy diagnostics, one campaign per task/method. The three metrics have different meanings and are not averaged.

Our harness records 41/51 reverse-identification hits, compared with 39/51 for direct and 36/51 for cosine. Sequential AUC3 improves by 0.004386 over direct, while the CEX learning arms have identical scores and executed-action hashes. The original Sequential confidence intervals compare policies with locked cosine, not harness with direct, and include zero for every learning arm. CEX's raw-PG reference is -0.007481; the shared learning-arm gain is 0.103608, with an original-scorer query-bootstrap 95\% interval of [0.068579, 0.139368]. This relative-objective improvement is shared by all three learning arms and is not harness-specific. Its absolute-response secondary objective worsens, so the result does not establish uniform drug-response benefit. Query-level replay intervals do not quantify variation across independent searches or biological laboratories.

The same-data federated variant does not outperform our non-federated harness. Across the three tasks, the latter uses 79 model calls and scores 17/18 candidates; the federated arm uses 87 calls and also scores 17/18; direct uses 18 calls and scores 9/18, with its other generations timing out at 900 seconds. These differences motivate the controlled DTI comparison, rather than an unsupported attribution of all score gains to research organization.

\subsection*{5.2 DTI held-out comparison}
\begin{center}\footnotesize\setlength{\tabcolsep}{3pt}
\begin{tabularx}{\linewidth}{>{\hsize=1.600\hsize\linewidth=\hsize\raggedright\arraybackslash}X>{\hsize=0.800\hsize\linewidth=\hsize\raggedright\arraybackslash}X>{\hsize=0.800\hsize\linewidth=\hsize\raggedright\arraybackslash}X>{\hsize=0.800\hsize\linewidth=\hsize\raggedright\arraybackslash}X}
\toprule
Method & Random AUROC & Unseen-drug AUROC & Unseen-protein AUROC \\
\midrule
DrugBAN & 0.8942 \(\pm\) 0.0230 & 0.8850 \(\pm\) 0.0317 & 0.8008 \(\pm\) 0.0734 \\
Released DrugBAN-Evo & 0.8715 \(\pm\) 0.0253 & 0.8612 \(\pm\) 0.0334 & 0.7181 \(\pm\) 0.0777 \\
Harness-free & 0.8950 \(\pm\) 0.0270 & 0.8853 \(\pm\) 0.0382 & 0.7958 \(\pm\) 0.0917 \\
Native DrugEvolve & 0.8963 \(\pm\) 0.0217 & 0.8865 \(\pm\) 0.0329 & 0.7975 \(\pm\) 0.0726 \\
Our harness & 0.9005 \(\pm\) 0.0207 & 0.8922 \(\pm\) 0.0299 & 0.8098 \(\pm\) 0.0678 \\
Cosine 5-NN (additional reference) & 0.8843 & 0.8719 & 0.7813 \\
\bottomrule
\end{tabularx}
\end{center}
Table 2. Neural rows use the fixed twelve-epoch refit, mean \(\pm\) sample SD over seeds 42--46; each searched row has five campaigns. All selected architectures completed refit without fallback and were sealed before testing. The additional deterministic cosine reference uses the identical 8,192 training pairs but was added retrospectively after the original tests; it has no training-seed SD.

Our harness has the highest observed mean AUROC on all three splits, above the fixed cosine reference. Primary differences are positive/negative/tied in 3/1/1 seeds against native DrugEvolve and 2/2/1 against direct generation. Both mean differences are below the predeclared 0.01 AUROC practical threshold, with intervals crossing zero (Table 3): a small, seed-dependent gain, not established superiority.

\begin{center}\footnotesize\setlength{\tabcolsep}{3pt}
\begin{tabularx}{\linewidth}{>{\hsize=1.600\hsize\linewidth=\hsize\raggedright\arraybackslash}X>{\hsize=0.850\hsize\linewidth=\hsize\raggedright\arraybackslash}X>{\hsize=0.850\hsize\linewidth=\hsize\raggedright\arraybackslash}X>{\hsize=0.850\hsize\linewidth=\hsize\raggedright\arraybackslash}X>{\hsize=0.850\hsize\linewidth=\hsize\raggedright\arraybackslash}X}
\toprule
Primary paired contrast & AUROC delta & 95\% bootstrap interval & Exact p & Holm p \\
\midrule
Our harness - Native DrugEvolve & +0.0057 & [-0.0174, +0.0280] & 0.6250 & 1.0000 \\
Our harness - Harness-free & +0.0069 & [-0.0147, +0.0326] & 0.7500 & 1.0000 \\
\bottomrule
\end{tabularx}
\end{center}
Table 3. Predeclared unseen-drug endpoint; paired campaign bootstrap and two-sided exact sign-flip tests (n=5).

Fixed-model variation reflects training seeds; searched-method variation also includes uncontrolled model generation. Secondary splits test transfer of selected architectures, not additional independent task families. Actual costs and fixed-order audit cases appear in Appendix D; Table D1 separates common ceilings from realized usage.

\subsection*{5.3 Matched proteomics results}
All 42 declared method/seed records are evaluated after both seals, with
35 unique prediction artifacts and no omitted seed or final-inference failure.
The final role contains 148 conditions, including 22 positives. Table 4
reports every searched arm and fixed reference; Appendix G gives all seed
values and secondary metrics.

\begin{table}[t]
% Generated from completed, independently rescored final results.
% results_sha256=e234c54386e4cda395d3273a992906595c436afce82799645709755b7139682c
\begin{center}\small\setlength{\tabcolsep}{7pt}
\begin{tabular}{lcc}
\toprule
Method & Final AP $\uparrow$ & Final AUROC $\uparrow$ \\
\midrule
Our harness (pooled) & $0.1822 \pm 0.0320$ & $0.5328 \pm 0.0488$ \\
Our harness + client cards & $0.1652 \pm 0.0157$ & $0.5034 \pm 0.0424$ \\
Harness-free (pooled) & $0.1926 \pm 0.0240$ & $0.5487 \pm 0.0371$ \\
Harness-free + client cards & $0.2002 \pm 0.0463$ & $0.5344 \pm 0.0774$ \\
Fixed efficacy head & $0.1589 \pm 0.0060$ & $0.4918 \pm 0.0224$ \\
Cosine 5-NN & $0.1685$ & $0.4989$ \\
Logistic regression & $0.1903$ & $0.6385$ \\
\bottomrule
\end{tabular}
\end{center}
Table N.4. Historical PTPC compound-disjoint final test, eight paired task-training seeds (mean $\pm$ sample SD). Cosine/logistic are deterministic references. Both pooled and cards arms access the same partner observations; cards add client-resolved diagnostics, including local AP/AUROC. The fixed head is our observed-response adaptation, not full ProteinTalks.

\end{table}

Harness-free+cards has the highest mean AP, 0.2002, compared with 0.1652
for our harness+cards. Their paired difference is $-0.0349$, with five
negative, two positive and one tied seed. Adding cards to our harness
also gives a negative mean difference, $-0.0170$. Neither predeclared
contrast supports the intended benefit (Table 5). The first descriptive
bootstrap interval is negative, but the exact Holm-adjusted test has
$p=0.125$; it does not establish multiplicity-corrected inferiority.

\begin{table}[t]
\begin{center}\footnotesize\setlength{\tabcolsep}{4pt}
\begin{tabular}{lcccc}
\toprule
AP contrast & Mean delta & 95\% interval & Exact p & Holm p \\
\midrule
Loop--direct (cards) & $-0.0349$ & $[-0.0623, -0.0101]$ & $0.0625$ & $0.1250$ \\
Cards--pooled (loop) & $-0.0170$ & $[-0.0429, +0.0028]$ & $0.4688$ & $0.4688$ \\
\bottomrule
\end{tabular}
\end{center}
Table N.5. Historical predeclared primary paired contrasts on AP. Intervals describe campaign variation conditional on this split; the practical threshold is 0.02 AP.

\end{table}

Absolute prediction quality is also limited. The neural methods' mean
AUROCs range from 0.4918 to 0.5487; logistic regression attains 0.6385
and BCE 0.4115, versus neural BCEs of 1.3719--1.4964. Logistic AP is
0.1903, above both loop variants. The final majority-class fraction is
126/148 (0.8514). Logistic attains exactly that accuracy by predicting every
condition negative at threshold 0.5; its ranking metrics use continuous
probabilities (Appendix J). Thus approximately 0.83 accuracy is not evidence
of strong efficacy discrimination. These results establish a completed observed-response
comparison with weak neural generalization, not an improvement over
ProteinTalks or a validated benefit from the federated diagnostic package.

\subsection*{5.4 Research cost and observed evidence transfer}
All 32 PTPC campaigns complete, with 128 successful fits, 640 client-worker receipts, 29 promotions and 67 rejections. Seven campaigns retain their starting head. Direct uses 24 calls per evidence mode; loop uses 85 with pooled diagnostics and 87 with cards. All input/output usage is observed. Loop therefore uses more calls and input tokens under the common ceilings; Tables H1--H2 report these costs separately from performance.

The fixed first-native cases give an executable example (Appendix H). In loop+pooled, seed 142, round 1, the proposal expects an explicit 24h-minus-6h representation to improve ranking. The code introduces shared per-time encoders, a difference branch and gated fusion. Visible dev/gate AP changes from 0.111316/0.630056 to 0.121785/0.752047, so the candidate is retained. Round 2 adds the drug embedding to the gate, expanding its input from 96 to 128; both AP values decrease, to 0.100287/0.727741, and the change is rejected. Round 3 inherits the retained source and explicitly cites the failed drug-conditioned gate before proposing head dropout. Source, decision, native-commit and next-context identities verify this sequence. It demonstrates recorded evidence transfer; held-out performance, rather than the generated explanation, determines method-level utility.
